\section{Introducción}

Con más de 600 millones de usuarios activos y un catálogo superior a los
100 millones de canciones, Spotify es la plataforma de streaming musical
más grande del mundo \cite{spotify2024stats}. El éxito de su sistema de
recomendación---responsable del 30\% de los plays---depende en gran medida
de técnicas de aprendizaje automático que identifican similitudes entre
canciones y preferencias de usuarios \cite{schedl2018current}.

El agrupamiento no supervisado ofrece una aproximación natural a la
organización de contenido musical: sin requerir etiquetas manuales,
permite descubrir estructuras latentes en el espacio de características
de audio. Variables como \textit{danceability}, \textit{energy} y
\textit{valence}, derivadas del análisis de señales de audio, capturan
aspectos perceptivos de la música que van más allá del género declarado
\cite{turnbull2008semantic}.

El presente trabajo aplica K-Means y análisis de componentes principales
(PCA) sobre el Spotify Tracks Dataset para identificar clusters musicales
coherentes. El objetivo práctico es construir la base algorítmica de un
sistema de recomendación que sugiera canciones similares basándose en
el cluster al que pertenecen.

\subsection{Objetivos}
\begin{itemize}
  \item Aplicar K-Means y métricas de evaluación (Silhouette, Davies-Bouldin)
        para determinar el número óptimo de clusters.
  \item Caracterizar cada cluster mediante el perfil promedio de sus
        características de audio.
  \item Visualizar los clusters en espacio reducido mediante PCA.
  \item Desplegar el sistema como aplicación web con Flask y Plotly.
\end{itemize}
